\documentclass[11pt,a4paper]{article}
\usepackage[T1]{fontenc}
\usepackage[utf8]{inputenc}
\usepackage{mathptmx}
\usepackage[margin=2.35cm]{geometry}
\usepackage{array,tabularx,longtable,pdflscape,amsmath,enumitem}
\usepackage{xcolor}
\usepackage[hidelinks]{hyperref}
\usepackage{fancyhdr}
\usepackage{microtype}
\definecolor{grid}{gray}{0.78}
\definecolor{headfill}{gray}{0.93}
\newcommand{\rowcolor}[1]{}
\newcolumntype{L}[1]{>{\raggedright\arraybackslash}p{#1}}
\setlength{\parindent}{0pt}
\setlength{\parskip}{6pt}
\setlength{\emergencystretch}{2em}
\renewcommand{\arraystretch}{1.22}
\setlist[itemize]{leftmargin=1.3em,itemsep=3pt,topsep=3pt}
\setlist[enumerate]{leftmargin=1.6em,itemsep=4pt,topsep=4pt}
\pagestyle{fancy}
\fancyhf{}
\fancyfoot[C]{\thepage}
\renewcommand{\headrulewidth}{0pt}
\renewcommand{\footrulewidth}{0pt}
\hypersetup{pdftitle={Usulan Penelitian RAG Salon dan Matriks RP RQ RO},pdfauthor={Matin Rusydan}}
\newcommand{\heading}[1]{\vspace{5pt}{\large\bfseries #1}\par\vspace{1pt}}
\newcommand{\subheading}[1]{\vspace{3pt}\textbf{#1}\par}

\begin{document}
\begin{center}
{\Large\bfseries REKOMENDASI USULAN PENELITIAN}\par
\vspace{4pt}
{\normalsize Matin Rusydan \quad | \quad 20 Oktober 2026}
\end{center}

\heading{Rencana Judul}
\textbf{Retrieval-Augmented Generation Berbasis BM25 dan Perutean Sumber\\untuk Konsultasi Layanan Salon}

\textbf{Topik:} Retrieval-Augmented Generation, perutean sumber, pencarian informasi heterogen, evaluasi jawaban konsultasi.

\heading{Masalah Utama}
Konsultasi layanan salon dapat membutuhkan dua jenis informasi dalam satu pertanyaan: fakta operasional yang berubah, seperti layanan dan harga, serta pengetahuan perawatan yang dijelaskan dalam dokumen. Penelitian pada pelanggan salon menempatkan layanan dan harga sebagai variabel yang berkaitan dengan kepuasan \cite{rachmawati}. Namun, temuan tersebut belum membuktikan adanya kesalahan informasi pada TIEN Salon. Masalah lokal perlu ditunjukkan melalui inventaris sumber, pertanyaan pengguna, dan pengujian awal sistem.

Retrieval-Augmented Generation (RAG) menyediakan konteks eksternal untuk membantu pembentukan jawaban. Akses database operasional melalui chatbot sudah diteliti oleh Burbano dkk. \cite{burbano}, sedangkan integrasi dokumen dan relasi pengetahuan pada layanan pelanggan dibahas oleh Hsueh dan Lin \cite{hsueh}. Karena itu, penggabungan database dan dokumen maupun penggunaan RAG pada aplikasi salon saja belum cukup menjadi kontribusi metode. Persoalan yang lebih spesifik adalah bagaimana memilih sumber yang berwenang untuk setiap fakta, sekaligus menjaga kelengkapan bukti pada pertanyaan campuran.

Sebagai contoh, pertanyaan ``Layanan untuk rambut kering yang tersedia apa, berapa harganya sekarang, dan bagaimana penjelasan perawatannya?'' memerlukan bukti dari database dan dokumen. Passage yang mirip secara semantik belum tentu memuat harga yang berlaku atau menjawab semua bagian pertanyaan. Penelitian layanan pelanggan mencatat kasus ketika sumber terstruktur yang lebih lama mengalahkan informasi terbaru \cite{hsueh}. Kajian evaluasi juga membedakan relevansi konteks, dukungan klaim, kebenaran jawaban, dan kualitas sitasi \cite{knollmeyer,gao}.

Penelitian ini mengusulkan \textbf{perutean sumber berdasarkan jenis fakta dan penyusunan konteks dengan prioritas sumber}: fakta operasional ditelusuri ke database yang berlaku, penjelasan perawatan ditelusuri ke dokumen terkurasi, dan pertanyaan campuran menggunakan keduanya. Pengaruh alokasi sumber dan pemisahan kueri diuji terhadap baseline BM25 pada korpus dan batas konteks yang sama. Implementasi BM25 berindeks dibandingkan dengan pemindaian ulang sebagai kontrol efisiensi dengan skor dan ranking yang harus setara. BM25 merupakan algoritma yang sudah ada; penggunaan algoritma tersebut bukan klaim kebaruan.

\textbf{Status usulan.} Rumusan ini merupakan arah eksperimen dari mapping 25 paper. Keberadaan kegagalan pada TIEN dan kebaruan spesifik mekanisme perutean masih harus dikonfirmasi melalui pilot dan penelusuran terarah. Judul ini bersifat sementara. Eksperimen komponen IR sudah tersedia pada development set, tetapi belum ada test set independen yang dikunci atau hasil evaluasi jawaban final. Tidak digunakan target skor minimal tanpa dasar kebutuhan dan bukti pengujian.

\clearpage
\begin{landscape}
\heading{Matriks Research Problems Research Questions dan Research Objectives}
Setiap baris menghubungkan satu masalah dengan pertanyaan yang dapat dijawab melalui eksperimen dan tujuan yang terukur. RQ1 berfokus pada fakta operasional, RQ2 pada kelengkapan bukti campuran, dan RQ3 pada kualitas keluaran serta biaya komputasinya.

\vspace{5pt}
\begin{tabularx}{\linewidth}{|>{\raggedright\arraybackslash}X|>{\raggedright\arraybackslash}X|>{\raggedright\arraybackslash}X|}
\hline
\rowcolor{headfill}\textbf{Research Problems (RP)} & \textbf{Research Questions (RQ)} & \textbf{Research Objectives (RO)}\\\hline
\textbf{RP1 -- Ketepatan sumber fakta operasional.}
Pencarian berdasarkan kemiripan teks dapat mengembalikan informasi layanan atau harga dari dokumen yang tidak berwenang atau tidak mutakhir. Integrasi database dengan RAG sudah tersedia dalam literatur, tetapi belum dengan sendirinya menjamin pemilihan sumber fakta yang benar pada kasus salon \cite{burbano,hsueh}. Kejadian dan besarnya kesalahan ini pada sistem lokal perlu diukur.
&
\textbf{RQ1.} Bagaimana perbedaan akurasi fakta layanan dan harga antara RAG dengan perutean sumber berbasis jenis fakta dan baseline BM25 pada snapshot data yang sama? Bagaimana ketepatan perilaku sistem ketika data aktif atau bukti promo/histori tidak tersedia?
&
\textbf{RO1.} Mengembangkan perutean fakta operasional ke database dan fakta perawatan ke dokumen, kemudian mengevaluasi akurasi serta kelengkapan field layanan/harga dibanding baseline BM25, dengan akses fakta dan snapshot yang setara. Uji ketidaktersediaan data dan invalidasi indeks ketika isi katalog berubah; jangan mengasumsikan tabel histori atau promo sudah tersedia.
\\\hline
\textbf{RP2 -- Kelengkapan bukti pertanyaan campuran.}
Satu pertanyaan dapat membutuhkan record layanan dan passage penjelasan perawatan sekaligus. Relevansi item individual atau ditemukannya dokumen asal belum menjamin seluruh kebutuhan jawaban mempunyai bukti. Evaluasi berbasis dokumen dapat melebihkan keberhasilan pencarian bukti \cite{topic,gao}.
&
\textbf{RQ2.} Bagaimana perbedaan cakupan bukti lengkap dan recall evidence pada pertanyaan campuran antara sistem usulan dan baseline, serta bagaimana perubahan hasil ketika pemisahan kueri dinonaktifkan dengan alokasi sumber tetap sama?
&
\textbf{RO2.} Menyusun konteks dari sumber yang diperlukan setiap pertanyaan dan mengukur cakupan bukti lengkap, Recall@k serta F1 retrieval. Melakukan ablation alokasi sumber tanpa pemisahan kueri untuk membedakan pengaruh kedua komponen. Penilaian cuplikan akhir dilakukan terpisah dari kecocokan ID evidence.
\\\hline
\textbf{RP3 -- Kualitas jawaban tidak terwakili oleh satu skor.}
Jawaban yang relevan belum tentu benar terhadap database atau didukung sitasi. Peningkatan metrik retrieval juga tidak selalu sejalan dengan faithfulness \cite{salman,es}. Evaluasi yang hanya melaporkan satu rata-rata dapat menutupi jawaban tidak lengkap, kesalahan sitasi, dan tambahan latensi \cite{knollmeyer,gao}.
&
\textbf{RQ3.} Bagaimana perubahan dukungan klaim, kebenaran dan kelengkapan jawaban, dukungan sitasi dokumen, serta waktu dan memori komponen retrieval pada sistem usulan dibanding baseline? Latensi end-to-end dianalisis terpisah ketika eksperimen generasi dijalankan.
&
\textbf{RO3.} Mengevaluasi keluaran dengan rubrik manual berbasis bukti, ukuran dukungan klaim dan sitasi, serta pencatatan biaya bangun indeks, waktu retrieval median/p95, dan alokasi memori indeks. Pisahkan pengukuran komponen dari biaya database, jaringan, dan generasi. Menganalisis trade-off dan kategori kegagalan secara berpasangan pada pertanyaan yang sama, tanpa mensyaratkan semua metrik meningkat bersamaan.
\\\hline
\end{tabularx}

\vspace{6pt}
\textbf{Batas klaim:} rumusan masalah menyatakan risiko yang didukung literatur dan perlu diverifikasi pada objek lokal. Ketiadaan suatu fitur pada satu paper tidak digunakan untuk menyatakan bahwa metode tersebut belum pernah diteliti secara global. Ablation utama menonaktifkan pemisahan kueri dengan akses sumber tetap setara. Percobaan tanpa salah satu sumber, jika ditambahkan, hanya merupakan diagnosis karena akses informasinya dikurangi.
\end{landscape}

\clearpage
\heading{Rancangan Metode dan Batas Penelitian}
\subheading{Sumber pengetahuan dan otoritas fakta}
\begin{itemize}
\item \textbf{Database salon:} layanan, harga, satuan, durasi dan syarat hanya jika field tersebut tersedia. Setiap evidence menyimpan ID record, nama field, versi snapshot dan waktu pembaruan.
\item \textbf{Dokumen perawatan:} dokumen yang telah ditelaah kesesuaian isi dan batas klaimnya. Simpan judul, sumber, tahun/tanggal pembaruan, halaman, bagian dan ID passage. Literatur algoritma RAG berfungsi sebagai referensi metode, bukan otomatis sebagai pengetahuan pelanggan.
\item \textbf{Web:} bila dibutuhkan, akuisisi mengikuti robots.txt dan aturan akses; log URL, waktu akses, cakupan dan kegagalan crawl. Konten tidak langsung dianggap benar hanya karena berhasil diambil.
\end{itemize}

\subheading{Alur sistem yang diusulkan}
Pertanyaan diidentifikasi kebutuhan faktanya sebagai operasional, dokumenter atau campuran. Sistem mengambil record database dan/atau passage dokumen sesuai kebutuhan, lalu membentuk konteks dengan identitas sumber. Fakta bisnis mengikuti database yang berlaku. Penjelasan perawatan harus ditautkan pada evidence dokumen. Jika bukti tidak tersedia atau pertanyaan ambigu, sistem menyatakan keterbatasan atau meminta klarifikasi berdasarkan aturan yang dicatat.

Perutean dimulai dengan aturan yang transparan dan dapat diaudit. Pilihan classifier, query decomposition berbasis model, atau komponen tambahan ditetapkan hanya jika pilot menunjukkan kebutuhan dan manfaatnya. Penelitian ini tidak memerlukan pelatihan model generatif baru.

\subheading{Pembanding dan pengujian komponen}
\begin{enumerate}
\item \textbf{B0 BM25 scan:} token, statistik korpus, dan skor dihitung ulang. Ini kontrol implementasi lama, bukan metode yang sengaja diberi informasi lebih sedikit.
\item \textbf{B1 BM25 berindeks:} statistik dan postings digunakan kembali. Skor serta ranking harus sama dengan B0; keuntungan yang diuji adalah efisiensi, bukan akurasi baru.
\item \textbf{A Alokasi tanpa pemisahan kueri:} slot database/dokumen dialokasikan untuk kebutuhan campuran, tetapi pencarian dokumen memakai pertanyaan asli.
\item \textbf{P Sistem usulan:} alokasi sumber disertai pemisahan kebutuhan fakta bisnis dan dokumen memakai aturan transparan. Parameter dipilih hanya pada development set.
\end{enumerate}

Semua metode memakai query, korpus, tokenisasi, parameter BM25 ($k_1=1{,}5$, $b=0{,}75$) dan $k$ yang sama. Uji cuplikan awal dan jendela berbasis kueri menggunakan batas non-web 400 karakter yang sama, tanpa melihat gold saat memilih cuplikan. Kualitas jawaban dibandingkan dengan generator, versi model dan anggaran konteks yang sama ketika generasi dijalankan. Dense/hybrid belum menjadi baseline yang dieksekusi dalam protokol ini.

Pada katalog aktif, isi fakta yang berubah harus memperbarui identitas indeks. Waktu pembacaan API saja tidak dianggap perubahan isi untuk ranking; provenance pembacaan tetap disimpan. Uji pembaruan harga harus memakai perubahan dan waktu observasi yang sama pada semua metode. Riwayat harga dan diskon tidak boleh diklaim tersedia sebelum skema/sumbernya benar-benar ada.

\subheading{Batas objek dan kontribusi}
Fokus kontribusi adalah ketepatan fakta, kelengkapan bukti, dan keterlacakan jawaban konsultasi. Penilaian diagnosis medis, dampak kepuasan pelanggan, serta akurasi klasifikasi citra tidak termasuk variabel penelitian ini. Kontribusi ilmiah ditentukan oleh perbedaan terukur terhadap baseline dan analisis kegagalan, bukan sekadar keberhasilan membuat website.

\clearpage
\heading{Definisi Parameter Uji}
Gold evidence ditetapkan pada tingkat record/field dan passage pendukung. Nomor dokumen saja tidak cukup untuk menyatakan bukti jawaban ditemukan. Untuk query campuran, daftar bukti dan fakta wajib dikelompokkan menurut subkebutuhan pertanyaan.

{\small
\begin{tabularx}{\linewidth}{|L{2.7cm}|X|L{4.0cm}|}
\hline
\rowcolor{headfill}\textbf{Ukuran} & \textbf{Objek dan cara penilaian} & \textbf{Hubungan dengan RQ}\\\hline
Akurasi fakta bisnis & Kecocokan layanan, harga, satuan dan syarat yang dinyatakan terhadap snapshot DB. Laporkan juga proporsi field wajib yang benar dan terisi agar penolakan tidak menaikkan skor secara semu. & Utama RQ1; pisahkan harga dari field lain.\\\hline
Precision@k, Recall@k, F1, MAP, MRR, nDCG & Relevansi item evidence yang diambil terhadap gold record/passage. TP, FP, FN dihitung pada unit yang sama. MAP menilai urutan semua hit; MRR memakai hit pertama. nDCG memakai label biner yang tersedia, bukan nilai bertingkat buatan. & RQ2; melaporkan mutu dan cakupan hasil pencarian.\\\hline
Cakupan bukti lengkap & Proporsi query yang seluruh kelompok gold evidence-nya ditemukan. Laporkan terpisah per kategori. Kecocokan ID harus dibedakan dari pemeriksaan apakah teks pendukung masih ada dalam cuplikan akhir. & Ukuran utama RQ2; definisi operasional proyek, bukan klaim metrik baru.\\\hline
Dukungan klaim & Klaim faktual pada jawaban yang didukung konteks yang benar-benar dipakai, dibagi klaim yang dinilai. & RQ3; faithfulness tidak sama dengan kebenaran harga terkini.\\\hline
Kebenaran dan kelengkapan & Rubrik manual untuk fakta wajib, ketepatan makna, dan informasi penting yang hilang terhadap gold answer bersumber. & RQ3; dibedakan dari relevansi jawaban.\\\hline
Dukungan sitasi & Kecocokan klaim dokumenter dengan passage yang dikutip serta cakupan klaim yang membutuhkan sitasi. Asal record DB dilaporkan terpisah. & RQ3; tautan sumber yang ada belum berarti mendukung klaim.\\\hline
Abstention dan coverage & Penolakan/klarifikasi yang benar dan keliru terhadap label answerability; proporsi pertanyaan yang dijawab. & Penjaga interpretasi RQ1--RQ3; cegah skor tinggi karena menolak semua query.\\\hline
Waktu dan memori & Biaya bangun indeks, median/p95 retrieval warm, serta alokasi Python pada indeks. End-to-end sejak query masuk hingga jawaban selesai diuji terpisah. & RQ3; catat perangkat, jumlah ulangan, cold/warm dan cakupan biaya.\\\hline
\end{tabularx}}

\vspace{6pt}
Untuk query $q$, dengan $E_q$ himpunan gold evidence dan $R_q^{(k)}$ hasil retrieval:
\[
 P_q@k=\frac{|R_q^{(k)}\cap E_q|}{k},\quad
 R_q@k=\frac{|R_q^{(k)}\cap E_q|}{|E_q|},\quad
 F1_q=\frac{2P_qR_q}{P_q+R_q}.
\]
Jika $P_q+R_q=0$, F1 ditetapkan 0. Jika hasil kosong, Precision dan F1 diberi 0 untuk query answerable; jika $E_q$ kosong, metrik retrieval tidak diterapkan dan perilaku abstention dievaluasi tersendiri. Precision@k memakai penyebut tetap $k$. Precision atas jumlah hasil aktual dilaporkan sebagai metrik berbeda. AP@k memakai jumlah seluruh gold evidence sebagai penyebut, termasuk gold di luar top-$k$; MAP adalah rata-rata AP antarkueri. RR memakai kebalikan peringkat hit pertama; MRR adalah rata-ratanya. nDCG memakai gain $2^{g}-1$ dan diskon $\log_2(r+1)$; label saat ini biner ($g=0$ atau $1$). Skor agregat utama adalah macro-average per query.

\textbf{F1 yang dimaksud di RQ2 adalah F1 retrieval}, bukan F1 klasifikasi citra, token-level answer F1, atau BERTScore-F1. Evaluasi manual digunakan agar tidak bergantung pada API judge berbayar. Jika RAGAS ditambahkan, versi, fungsi metrik dan model evaluator harus tercatat \cite{es}; rubrik buatan sendiri tidak diberi label implementasi RAGAS.

\clearpage
\heading{Tugas Berikutnya Setelah Mapping}
\textbf{Tahap terdekat adalah validasi masalah dan pilot, kemudian penetapan metode.} Pengembangan program difokuskan pada komponen yang dibutuhkan untuk menghasilkan perbandingan yang adil.

\begin{enumerate}
\item \textbf{Konfirmasi celah spesifik.} Telusuri studi yang paling dekat tentang perutean sumber, structured--unstructured retrieval, source authority dan temporal consistency. Buat matriks pembeda: jenis sumber, mekanisme routing, konflik informasi, kategori query dan evaluasi. Integrasi database/dokumen tidak diklaim baru.
\item \textbf{Buktikan masalah lokal.} Inventaris field layanan, asal pengetahuan dan jalur jawaban sistem. Kumpulkan contoh pertanyaan operasional, dokumenter dan campuran. Pisahkan pertanyaan pelanggan nyata dari pertanyaan yang disusun peneliti.
\item \textbf{Bangun pilot berlabel.} Pilot development saat ini berisi 24 query dalam kategori fakta DB, dokumen, campuran, dan kasus sulit. Dua puluh query memiliki gold evidence retrieval; empat lainnya memerlukan penilaian kebijakan atau keterjawaban. Kasus sulit mencakup sumber konflik, perubahan harga, ambigu atau tidak terjawab. Jumlah ini untuk diagnosis awal, bukan kecukupan statistik penelitian final.
\item \textbf{Uji baseline dan analisis error.} Jalankan BM25 scan, BM25 berindeks, alokasi tanpa pemisahan kueri, dan sistem lengkap. Simpan evidence aktual, cuplikan, versi konfigurasi, status error dan waktu; jawaban hanya dicatat bila generasi benar-benar dijalankan. Temukan apakah kegagalan dominan berasal dari routing, retrieval, data, atau generation. Pastikan evaluator tidak melakukan retrieval ulang yang berbeda dari konteks generasi.
\item \textbf{Kunci intervensi dan data final.} Jika pilot menunjukkan masalah relevan, tetapkan mekanisme perutean serta prioritas sumber. Susun development set terpisah dan test set yang dikunci; tentukan jumlah dan strata final berdasarkan keragaman kasus serta ketidakpastian yang ingin diukur.
\item \textbf{Lakukan eksperimen berpasangan.} Bandingkan baseline, sistem usulan dan ablation pada query sama. Laporkan hasil per kategori, trade-off, contoh kegagalan, dan ketidakpastian sesuai data. Finalisasi judul, abstrak dan simpulan setelah hasil tersedia.
\end{enumerate}

\subheading{Status eksekusi 20 Oktober 2026}
Eksperimen IR offline pada 13 evidence dan 24 query development telah dijalankan. Pada 20 query ber-gold, baseline BM25 memperoleh Recall@5 $0{,}975$ dan cakupan kelompok bukti lengkap $19/20$; alokasi disertai pemisahan kueri memperoleh $1{,}000$ dan $20/20$. Precision@5 ketat masing-masing $0{,}240$ dan $0{,}250$. BM25 scan dan berindeks menghasilkan ranking yang sama. Ini hasil diagnosis development dengan anotasi revisi yang masih menunggu peninjauan, bukan bukti generalisasi atau kualitas jawaban generatif. Tidak ada panggilan LLM pada eksperimen ini.

Pengukuran waktu menggunakan korpus lokal yang dimuat sebelumnya dan diperbanyak secara sintetis; hasilnya tidak mewakili cold start Vercel, concurrency, keseluruhan memori aplikasi, atau biaya generasi. Semua jawaban dan dukungan sitasi yang belum diperiksa tetap berstatus belum dinilai. Template pemeriksaan acuan, span cuplikan, dan keluarga skenario disediakan; test set final belum dikunci.

\subheading{Artefak yang perlu tersedia sebelum eksperimen final}
Inventaris sumber beserta otoritas field; snapshot database; korpus dokumen terkurasi; 24 query pilot dengan gold answer dan evidence; rubrik anotasi; baseline yang dapat direproduksi; dan laporan error pilot. Itulah dasar untuk memutuskan apakah metode usulan layak diteruskan atau perlu disederhanakan.

\subheading{Aturan validitas evaluasi}
Skor nol yang sah tetap dihitung. Error API/parsing dicatat terpisah dengan denominator transparan; laporkan tingkat keberhasilan eksekusi dan hasil end-to-end. Jangan memilih hanya pertanyaan dengan seluruh metrik non-zero. Anotasi satu orang dinyatakan sebagai keterbatasan; jika pemeriksa kedua tersedia, lakukan verifikasi sebagian dengan protokol yang dicatat. Hindari klaim inter-annotator agreement bila tidak ada dua annotator.

\subheading{Gate keputusan pilot}
Jika baseline sudah benar dan lengkap pada sebagian besar kasus yang relevan, periksa apakah dataset terlalu mudah atau masalah yang diusulkan tidak cukup penting. Jika perbaikan tidak terbukti, laporkan hasil tersebut dan ubah fokus berdasarkan error yang nyata. Jangan memaksakan target akurasi tertentu atau mengubah test set agar metode tampak unggul.

\clearpage
\renewcommand{\refname}{Referensi Utama}
\begin{thebibliography}{9}
\setlength{\itemsep}{8pt}
\bibitem{rachmawati} I. Rachmawati dkk. ``Pengaruh kualitas layanan, kualitas produk dan harga terhadap kepuasan pelanggan salon di Tanjung Tabalong Kalimantan Selatan.'' \textit{Jurnal Akuntansi, Ekonomi dan Manajemen Bisnis}, 5(3), 588--597, 2025. DOI: \href{https://doi.org/10.55606/jaemb.v5i3.7103}{10.55606/jaemb.v5i3.7103}. Mapping P20.
\bibitem{burbano} J. Burbano, P. Landeta-L\'opez, C. Guevara-Vega, dan A. Qui\~na-Mera. ``Design and Evaluation of a Retrieval-Augmented Generation LLM Chatbot with Structured Database Access.'' \textit{Applied Sciences}, 16, 3147, 2026. DOI: \href{https://doi.org/10.3390/app16073147}{10.3390/app16073147}. Mapping P01. Lihat catatan kritis statistik/waktu pada mapping; angka peningkatannya tidak digunakan sebagai target usulan.
\bibitem{hsueh} N.-L. Hsueh dan W.-C. Lin. ``Design of a RAG-Based Customer Service Chatbot Enhanced with Knowledge Graph and GPT Evaluation: A Case Study in the Import Trade Industry.'' \textit{Software}, 5, 15, 2026. DOI: \href{https://doi.org/10.3390/software5020015}{10.3390/software5020015}. Mapping P02.
\bibitem{knollmeyer} S. Knollmeyer dkk. ``Evaluating Retrieval Augmented Generation: A Comprehensive Review of Evaluation Dimensions, Question Types, and Application.'' \textit{SN Computer Science}, 7, 576, 2026. DOI: \href{https://doi.org/10.1007/s42979-026-05134-x}{10.1007/s42979-026-05134-x}. Mapping P09.
\bibitem{gao} T. Gao, H. Yen, J. Yu, dan D. Chen. ``Enabling Large Language Models to Generate Text with Citations.'' \textit{Proceedings of EMNLP}, 6465--6488, 2023. \url{https://aclanthology.org/2023.emnlp-main.398/}. Mapping P11; prosiding terbit.
\bibitem{salman} M. D. Salman, Rahmaddeni, T. Nasution, dan Susanti. ``Implementation of Retrieval-Augmented Generation Method on Large Language Model for Development of Campus Service and Information Chatbot.'' \textit{INOVTEK Polbeng -- Seri Informatika}, 11(1), 298--309, 2026. DOI: \href{https://doi.org/10.35314/3y9hy151}{10.35314/3y9hy151}. Mapping P05.
\bibitem{topic} K. Topi\'c, M. Bagi\'c Babac, dan V. Mornar. ``Evaluating LLM-Based Retrieval-Augmented Generation for Soil Science Question Answering.'' \textit{Information}, 17, 859, 2026. DOI: \href{https://doi.org/10.3390/info17090859}{10.3390/info17090859}. Mapping P07.
\bibitem{es} S. Es, J. James, L. Espinosa-Anke, dan S. Schockaert. ``RAGAS: Automated Evaluation of Retrieval Augmented Generation.'' \textit{Proceedings of EACL: System Demonstrations}, 150--158, 2024. \url{https://aclanthology.org/2024.eacl-demo.16/}. Mapping P10; prosiding terbit.
\end{thebibliography}
\end{document}
